\scanpage{93}
\mspage{10}
% archive sheet number: 92

2. \uline{Comparaison de la cohomologie de la fibre spéciale et générale par la suite spectrale de Leray.}

2.1. $X$ propre sur $S$, $U$ ouvert de $X_{\unsure{\eta|y}}$.

La suite spectrale de \deleted{localisation} \inserted{cohomologie relative} sera ici remplacée, essentiellement, par la suite spectrale \inserted{(légèrement différente)} de l'inclusion ouverte $U_\alpha\longrightarrow X_\alpha$, ou plutôt par la limite de ses suites spectrales. \deleted{\unsure{La plus}} Nous ne supposons pas ici que $k(s)$ soit séparablement clos, car nous désirons \deleted{\unsure{générale}} \inserted{parfois} suivre les opérations de $\unsure{\pi|\Pi}$ tout entier, pas de $I$ seulement. Cela nous amène à introduire, en plus de l'ouvert $U$ de $X_\eta$ et de $\bar U$, $\bar X=X_{\bar\eta}$, les schémas \unsure{analogues|conjugués} $\tilde X$ et $\tilde U$ sur $\tilde S$.

\diagram{\begin{tikzcd}[ampersand replacement=\&]
U \arrow[d,phantom,"\cap" description] \& \tilde U \arrow[d,phantom,"\cap" description] \& \bar U=U_{\bar\eta} \arrow[d,phantom,"\cap" description] \\
X \arrow[d,no head] \& \tilde X \arrow[l] \arrow[d,no head] \& \bar X=X_{\bar\eta} \arrow[l,"\unsure{g|\beta}"] \arrow[d,no head] \\
S \& \tilde S \arrow[l] \& \bar S \arrow[l,no head]
\end{tikzcd}\quad (2,2)}

\deleted{On utilise la suite} Le groupe $\pi$ opère sur la situation
\diagram{\begin{tikzcd}[ampersand replacement=\&]
\tilde U \arrow[d,phantom,"\cap" description] \& \bar U \arrow[l] \arrow[d,phantom,"\cap" description] \\
\tilde X \& \bar X \arrow[l]
\end{tikzcd}}
le sous-groupe $I$ opérant par $\tilde X$-automorphismes.

Nous considérons la suite spectrale de Leray pour le composé \deleted{$\bar U$}
\[
\tilde X\xleftarrow{g}\bar U\ ,
\tag*{(2,3).}
\]
on trouve
\[
H^*(\bar U,\Lambda_\nu)\Longleftarrow E_2^{pq}=H^p(\tilde X,{\unsure{R|\mathbb R}}^qg_*((\Lambda_\nu)_{\bar U})\ .
\tag*{\deleted{\unsure{(2.3)}} (2.4)}
\]
\deleted{On peut considérer cette suite spectrale comme une suite}
Le th de changement de base propre permet de l'écrire sous la forme
